\documentclass[11pt]{article}
\usepackage{polyglossia}
\setdefaultlanguage{romanian}
\author{M.~Gătejescu}
\title{Zei Americani de Neil Gaiman}
\begin{document}

\maketitle

Pentru mult timp am evitat cartea, crezând că este mai mult o carte
celebră\slash populară, decât o carte de substanță. Am devenit
interesat după ce am aflat că Zei Americani a câștigat un număr
mare de premi, printre care Hugo și Nebula. După ce am citit-o
pot să spun că nu am fost dezamăgit.

Așa cum spune titlul romanul abordeză problema mitologies, dintr-un
punct de vedere forte modern și foarte actual.

Deci tema principală este lupta dintre zei, zeii vechi, printe aceștia
se numără Odin, Loki și Lebda, și zeii noi care sunt întruchipări
ale ideilor actuale care acaparează conștiința colectivă, cum ar
fi internetul, băncile, televiziunea, teorile conspiraționiste etc.

Nu pot să nu obserb reluare motivului lui Iisus și a sacrificilui,
în Shadow. Care se sacrifică pentru Odin, tatăl lui, călătoreșe în
lumea de dincolo și este adus înapoi. Acțiunile lui următoare opresc
războiul zeilor și vindecă\slash regenerează relațiile dintre zei,
lumea ca întreg. Purificare continuă și crimele din Lakeville sunt
oprite.
Acest sfâșit mă lasă cu multe întrebări: Dacă Odin s-a renăscut,
va fi și Loki metamorfozat? Nu pare plauzibil\ldots Ce se întâmplă
cu noi zei, eu sunt o personificare a multor defecte are societății
actuale, se vor shimba și ei.
Cartea nu răspunde la aceaste întrebări generale, la nivel macro,
poate pentru că nu se știu răspunsurile. Autorul se adresează
cititorului, individului, și îi atrage atenția că deși trăim
într-o epocă a digitalului și care până recent părea dominată
de rațional, zeii încă există, atât noi, cât și vechi.

\section*{}
Notă legată de traducere: Mi-a plăcut că numele, mai ales numele
personajului principal, Shadow, nu au fost traduse.

Notă: Neil Gaiman este un scriitor prolific. Probabil că o să mai
citesc cărți scrise de el. Am adăugat Băieții lui Anansi pe lunga
listă de lectură.

Traducere din limba română și note de Liviu Radu

ISBN 978-606-93637-7-5

\end{document}
